\subsection{Higher moments in the minimal singularity type}
\label{sec:big-minimal-moments}

In this subsection the reference is the potential $V=V_\theta=P_\theta(0)$ with minimal singularities. It is an ordinary model, and we write $\mathcal E_V=\mathcal E_{\mathrm{env}}(X,\theta,V)$ and
\[
 \ell_V=\underline{\operatorname{vol}}_\theta(V)\geq\ell>0,
 \qquad M_V=\overline{\operatorname{vol}}_\theta(V)<\infty.
\]
The first inequality follows from volume monotonicity, since $\varphi\preceq V$, and the last from bounded mass. Thus $V$ satisfies the standing assumptions on $\varphi$, and the results above apply with $\varphi$ replaced by $V$. For a probability measure $\nu$, $A_\theta(\nu)$ denotes the supremum of $\int_X(-w)\,\mathrm{d}\nu$ over the $w\in\PSH(X,\theta)$ with $\sup_Xw=0$, as in \cref{subsec:psh-measure-smoothing}. We show that a first moment of $V-u$ against the measure $\mu_u$ of $u$ itself controls a higher moment of $V-u$ against every fixed $\nu$ with $A_\theta(\nu)<\infty$. An envelope with a concave obstacle converts the distribution of $V-u$ under $\mu_u$ into this stronger integrability estimate. We state the envelope lemma for the reference $\varphi$; here it is applied with $V$ in place of $\varphi$, and in \cref{sec:big-model-lp} with $\varphi$ itself.

\Needspace{20\baselineskip}
\begin{lemma}\label{lem:big-concave-envelope}\label{lem:big-model-concave-envelope}
Let $u\in\mathcal E_\varphi$ with $u\leq\varphi$, and let
$\chi\colon\mathbb R_{\leq0}\to\mathbb R_{\leq0}$ be increasing and concave, with $\chi(0)=0$. Suppose that $\chi$ is $C^{1,1}$ on compact intervals and that, for a finite constant $L\geq1$,
\[
 1\leq\chi'\leq L.
\]
Then the envelope
\[
 p=P_\theta(\varphi+\chi(u-\varphi))
\]
belongs to $\mathcal E_\varphi$, and
\[
 \mu_p\leq\mathbf1_{\{p=\varphi+\chi(u-\varphi)\}}
                 \chi'(u-\varphi)^n\mu_u.
\]
Here the obstacle and the contact set are defined off the pluripolar set where $u-\varphi$ is undefined, which all these measures ignore.
\end{lemma}
\begin{proof}
For $t\leq0$, the bounds on $\chi'$ give $Lt\leq\chi(t)\leq t$. Thus
\[
 Q_L(u,\varphi)\leq p\leq u\leq\varphi.
\]
The lower potential belongs to $\mathcal E_\varphi$ by \cref{lem:relative-stretch-closure}, so $p$ is not identically $-\infty$ and belongs to $\mathcal E_\varphi$ by upward stability. The obstacle is finite and quasi-continuous off the poles, so the contact theorem \cref{thm:big-general-envelope-contact} applies.

Let $\tau=\chi^{-1}$. It is convex, with $L^{-1}\leq\tau'\leq1$. We justify the transformation
\[
 w=\varphi+\tau(p-\varphi)
\]
without assuming that $\varphi$ is locally bounded on some open set. On a coordinate ball choose a smooth function $q$ with $\theta+\ddc q\geq0$. The functions
\[
 p_m=\max(p,q-m),\qquad
 \varphi_m=\max(\varphi,q-m)
\]
are locally bounded $\theta$-plurisubharmonic functions on the ball, and $p_m\leq\varphi_m$. \Cref{lem:bounded-convex-pair} applies to
$w_m=\varphi_m+\tau(p_m-\varphi_m)$ and gives
\[
 \theta_{w_m}\geq0,\qquad
 \mu_{w_m}\geq\tau'(p_m-\varphi_m)^n\mu_{p_m}.
\]

The function $(s,t)\mapsto t+\tau(s-t)$ on $\{s\leq t\}$ is nondecreasing in each variable. Consequently $w_m$ decreases as $m$ increases. Since $p\leq p_m\leq w_m\leq\varphi_m$, the limit is not identically $-\infty$ on the ball. Off the poles it equals $w$. The local limits therefore agree on overlaps after regularization and define a global $\theta$-plurisubharmonic function.

On $D_m=\{p>q-m\}$, neither $p$ nor $\varphi$ is truncated. Thus $w_m=w$ there, and plurifine locality transfers the product inequality to $D_m$. These sets exhaust the complement of the poles of $p$ in the ball. Since all products ignore the discarded pluripolar set,
\[
 \mu_w\geq\tau'(p-\varphi)^n\mu_p.
\]
The obstacle bound $p\leq\varphi+\chi(u-\varphi)$ gives $w\leq u$, with equality on the contact set of $p$. Ordered contact, the identity
$\tau'(p-\varphi)=1/\chi'(u-\varphi)$ on that set, and contact support of $\mu_p$ prove the assertion.
\end{proof}

\Needspace{11\baselineskip}
\begin{proposition}\label{prop:big-minimal-higher-moment}
Set $N=2n+1$ and $r=1+(4N)^{-1}$. For $u\in\PSH(X,\theta)$ with $u\sim V$ and $\sup_Xu=0$, put $B=\int_X(V-u)\,\mu_u$. There is a constant $C$, depending only on $n,\ell_V,M_V$, such that, for every non-pluripolar probability measure $\nu$ with $A=A_\theta(\nu)<\infty$,
\[
 \int_X(V-u)^r\,\mathrm{d}\nu
 \leq C(1+A)(1+B)^{1+1/N}
 .
\]
In particular, normalized potentials $u\sim V$ with uniformly bounded $B$ are uniformly integrable with respect to $\nu$.
\end{proposition}
\begin{proof}
From the distribution of $V-u$ under $\mu_u$ we build a concave weight adapted to $u$, apply the preceding lemma, and use $A$ to bound the resulting envelope against $\nu$. Throughout, $c$ and $C$ denote positive constants depending only on $n,\ell_V,M_V$.

The normalization gives $u\leq V$. Put $x=V-u\geq0$ and
\[
 T=\sup\{t\geq0:\mu_u(\{x>t\})>0\},
\]
with $T=0$ if this set is empty. Since $u\sim V$, $T$ is finite. Zero domination, applied to $u$ and $V-t$ for $t>T$, shows that $u\geq V-T$; thus $x\leq T$ off the pole set. If $T\leq2$ the claim is immediate. We may therefore suppose $T>2$; in particular $B>0$.

Fix a real number $s$ with $0<s<T-2$. For $t\geq0$, define
\[
 g_s(t)=
 \begin{cases}
  \bigl((1+t)^{3/2}\mu_u(\{x>t\})\bigr)^{-1},&0\leq t\leq T-s,\\
  (1+t)^{-3/2},&t>T-s,
 \end{cases}
 \quad F_s(t)=1+\int_0^t g_s(a)\,\mathrm{d}a,
\]
and set
\[
 W_s(t)=\int_0^t F_s(a)^{1/N}\,\mathrm{d}a,
 \qquad\chi_s(-t)=-W_s(t).
\]
The tail $\mu_u(\{x>t\})$ in the first denominator has a positive lower bound on $[0,T-s]$. Hence $F_s$ is bounded, nondecreasing and locally Lipschitz, with $F_s\geq1$. Since $u\leq V$ and, by \cref{prop:relative-elementary}, $u\in\mathcal E_V$, the function $\chi_s$ satisfies the hypotheses of \cref{lem:big-concave-envelope}, with $V$ in place of $\varphi$.

Tonelli's theorem gives
\begin{equation}\label{eq:big-weight-derivative-moment}
 \int_X F_s(x)\,\mu_u
 =\mu_u(X)+\int_0^Tg_s(t)\mu_u(\{x>t\})\,\mathrm{d}t
 \leq3\max(1,M_V),
\end{equation}
since $\int_0^\infty(1+t)^{-3/2}\,\mathrm{d}t=2$ and the tail is at most $M_V$ on the second interval.
Chebyshev's inequality gives $\mu_u(\{x>t\})\leq B/t$. Integrating first for $F_s$ and then for $W_s$, we obtain
\begin{equation}\label{eq:big-weight-growth}
 W_s(t)\geq c(1+B)^{-1/N}t^{\beta},
 \qquad1\leq t\leq T-s,
 \qquad\beta=1+\frac1{2N}.
\end{equation}
Indeed, $g_s(t)\geq cB^{-1}t^{-1/2}$ for $t\geq1$ in this interval, so $F_s(t)\geq c(1+B)^{-1}t^{1/2}$. For $t\geq2$ integrate this estimate over $[t/2,t]$; for $1\leq t\leq2$ use $W_s(t)\geq t$.

Let $p_s=P_\theta(V-W_s(x))\in\mathcal E_V$, the envelope of \cref{lem:big-concave-envelope} for $\chi=\chi_s$. The lemma, the bound $W_s(t)\leq tF_s(t)^{1/N}$ and the Cauchy--Schwarz inequality give, since $N=2n+1$,
\[
 \begin{split}
 \int_X(V-p_s)^{1/2}\,\mu_{p_s}
 &\leq\int_X x^{1/2}F_s(x)^{(n+1/2)/N}\,\mu_u\\
 &\leq B^{1/2}\left(\int_XF_s(x)\,\mu_u\right)^{1/2}
 \leq C B^{1/2}.
 \end{split}
\]
Write $a_s=\sup_Xp_s\leq0$. By the definition of $V$, $p_s-a_s\leq V$, so $V-p_s\geq-a_s$. Since $\mu_{p_s}(X)\geq\ell_V$, the last estimate yields $\ell_V\sqrt{-a_s}\leq C\sqrt B$, and hence $-a_s\leq CB$.

The definition of $A$ now gives
\[
 \int_XW_s(x)\,\mathrm{d}\nu
 \leq\int_X(V-p_s)\,\mathrm{d}\nu
 \leq A-a_s\leq A+CB.
\]
Combining this with \eqref{eq:big-weight-growth}, and then letting $s\downarrow0$ for each fixed $1\leq t<T$, we obtain
\[
 \nu(\{x>t\})\leq C(1+B)^{1/N}(A+B)t^{-\beta}.
\]
There is no tail for $t\geq T$. The layer-cake formula therefore gives
\[
 \int_Xx^r\,\mathrm{d}\nu
 =r\int_0^\infty t^{r-1}\nu(\{x>t\})\,\mathrm{d}t
 \leq1+C(1+B)^{1/N}(A+B),
\]
because $1<r<\beta$. This implies the stated bound.

Finally $-u=-V+(V-u)$, where $-V\in L^1(\nu)$ because $A<\infty$. The first summand is fixed and integrable, and the second has a uniform $L^r(\nu)$ bound. This proves uniform integrability; no $L^r$ bound for $V$ itself is claimed.
\end{proof}
